% !TeX root = ../../Book5.tex
% 纲要标题：### 5.3 Brauer 诱导
% 纲要位置：计划纲要.md，第 3214 行。

\section{Brauer 诱导}
\label{b5:sec:0503}

% TODO: 按以下纲要要点撰写本节。

% 纲要内容（仅作写作计划，尚非教材正文）：
%
% * 初等子群及一维特征标
% * Brauer 诱导定理的明确版本
% * 与原纲要未指明版本的“Brauer 定理”对接
% * 定义 \(p\)-初等群为 \(P\times C\)，其中 \(P\) 为有限 \(p\)-群、\(C\) 为阶与 \(p\) 互素的循环群；所有素数一起使用
% * 目标版本为复虚特征标是从初等子群诱导的特征标的整数线性组合，并细化到适当子群的一维特征标；完整整性证明在第5.4节
% * 本章不以前置第8章的模表示、Brauer 特征标或分解数为条件
%

\subsection{初等子群与一维特征标}
\label{b5:subsec:0503:01}

待撰写。

\subsection{Brauer 诱导定理的整系数版本}
\label{b5:subsec:0503:02}

待撰写。

\subsection{Brauer 诱导定理的版本辨析}
\label{b5:subsec:0503:03}

待撰写。

\subsection{\texorpdfstring{\(p\)-初等群的定义}{p-初等群的定义}}
\label{b5:subsec:0503:04}

待撰写。

\subsection{从初等子群到一维特征标的诱导}
\label{b5:subsec:0503:05}

待撰写。

\subsection{Brauer 诱导与模表示的关系}
\label{b5:subsec:0503:06}

待撰写。
